% =====================================================================
% Speaker notes for the defence presentation — one page per slide.
% English (short read-aloud sentences) + Arabic (for rehearsal).
% Compile with XeLaTeX.
% =====================================================================
\documentclass[12pt,a4paper]{article}
\usepackage[left=2.2cm,top=2.0cm,right=2.2cm,bottom=2.0cm]{geometry}
\usepackage{newtxmath}
\usepackage{fontspec}
\usepackage{polyglossia}
\setmainlanguage{english}
\setotherlanguage{arabic}
\newfontfamily\arabicfont[Script=Arabic,Scale=1.1]{Traditional Arabic}
\usepackage{xcolor}
\definecolor{honest}{HTML}{0072B2}
\linespread{1.15}
\setlength{\parindent}{0pt}
\setlength{\parskip}{5pt}
\pagestyle{plain}

\newcommand{\slidehead}[2]{%
\noindent\colorbox{honest!10}{\parbox{\dimexpr\textwidth-2\fboxsep}{%
\textbf{#1}\hfill\textbf{#2}}}\par\vspace{4pt}}
\newcommand{\arstart}{\par\vspace{6pt}\hrule\vspace{6pt}}

\begin{document}

% ---------------- Slide 1 ----------------
\slidehead{Slide 1 --- Title}{target 0:30}
Good morning. My name is Suhail Alhammali.
My project is about the voice and Parkinson's disease.
I will show how a computer can measure a voice recording.
And I will show how we tested it honestly.
The word ``honestly'' is the heart of this talk.
\arstart
\begin{Arabic}
صباح الخير. اسمي سهيل الحمالي.
مشروعي عن الصوت ومرض باركنسون.
سأعرض كيف يستطيع الحاسوب أن يقيس تسجيلاً صوتياً.
وسأعرض كيف اختبرنا النظام بأمانة.
وكلمة «بأمانة» هي قلب هذا العرض.
\end{Arabic}
\clearpage

% ---------------- Slide 2 ----------------
\slidehead{Slide 2 --- A movement disease reaches the voice}{target 1:00}
Parkinson's disease destroys dopamine cells in the brain.
These cells control movement.
So patients get tremor, stiffness, and slow movement.
But the disease weakens every muscle.
That includes breathing muscles, the larynx, the tongue, and the lips.
Follow the chain on the slide.
Dopamine loss leads to poor motor control.
That causes a speech disorder called hypokinetic dysarthria.
And that produces acoustic changes we can measure.
\arstart
\begin{Arabic}
يُتلف مرض باركنسون خلايا الدوبامين في الدماغ.
وهذه الخلايا تتحكم في الحركة.
لذلك يصاب المرضى بالرعاش والتيبّس وبطء الحركة.
لكن المرض يُضعف كل عضلة في الجسم.
ويشمل ذلك عضلات التنفّس والحنجرة واللسان والشفتين.
تابعوا السلسلة على الشريحة.
فقدان الدوبامين يؤدي إلى ضعف التحكم الحركي.
وهذا يسبّب اضطراب كلام اسمه عُسر التلفّظ ناقص الحركة.
وذلك يُنتج تغيّرات صوتية يمكننا قياسها.
\end{Arabic}
\clearpage

% ---------------- Slide 3 ----------------
\slidehead{Slide 3 --- The voice is a cheap, early window}{target 0:55}
Why the voice? Three reasons.
First, about ninety percent of patients show speech changes.
Second, the changes can start early.
In one documented case, about five years before diagnosis.
Third, a microphone is cheap and safe.
No needles. No scans. You can repeat it any time.
To be clear: this is not a diagnosis.
It is a reason to see a doctor earlier.
\arstart
\begin{Arabic}
لماذا الصوت؟ لثلاثة أسباب.
أولاً، نحو تسعين بالمئة من المرضى تظهر عندهم تغيّرات في الكلام.
ثانياً، قد تبدأ التغيّرات مبكراً.
في حالة موثّقة واحدة، قبل التشخيص بنحو خمس سنوات.
ثالثاً، الميكروفون رخيص وآمن.
لا إبر ولا أشعة. ويمكن تكراره في أي وقت.
وللتوضيح: هذا ليس تشخيصاً.
بل سبب لزيارة الطبيب مبكراً.
\end{Arabic}
\clearpage

% ---------------- Slide 4 ----------------
\slidehead{Slide 4 --- The research question}{target 0:30}
Here is the whole project in one question.
Can a computer measure a voice recording?
And can it say which group the pattern resembles?
Patients with Parkinson's, or healthy speakers?
Everything that follows answers this question.
\arstart
\begin{Arabic}
هذا هو المشروع كله في سؤال واحد.
هل يستطيع الحاسوب أن يقيس تسجيلاً صوتياً؟
وهل يستطيع أن يقول أي مجموعة يشبه النمط؟
مرضى باركنسون أم الأشخاص الأصحّاء؟
وكل ما يأتي بعد ذلك يجيب عن هذا السؤال.
\end{Arabic}
\clearpage

% ---------------- Slide 5 ----------------
\slidehead{Slide 5 --- Six stages, one shared implementation}{target 1:15}
This diagram is the spine of the system.
You will see it again on later slides.
Raw audio comes in on the left.
We clean it in five fixed steps.
We cut it into ten-second chunks.
From each chunk we take seventy-four measurements.
A model turns them into a score.
The output is a cautious indication, not a verdict.
One important point: the same code runs everywhere.
Training, testing, and prediction use one implementation.
The program checks this before every prediction.
\arstart
\begin{Arabic}
هذا المخطط هو العمود الفقري للنظام.
وسترونه مرة أخرى في شرائح لاحقة.
يدخل الصوت الخام من اليسار.
ننظّفه في خمس خطوات ثابتة.
ثم نقطّعه إلى مقاطع مدتها عشر ثوانٍ.
ومن كل مقطع نأخذ أربعة وسبعين قياساً.
ثم يحوّلها نموذج إلى درجة.
والمخرج مؤشر حذر، وليس حكماً.
ونقطة مهمة: الشيفرة نفسها تعمل في كل مكان.
فالتدريب والاختبار والتنبّؤ تستخدم تنفيذاً واحداً.
والبرنامج يتحقق من ذلك قبل كل تنبّؤ.
\end{Arabic}
\clearpage

% ---------------- Slide 6 ----------------
\slidehead{Slide 6 --- Every measurement points at physiology}{target 1:15}
We did not choose measurements at random.
Each one reflects something physical in the throat or mouth.
Monotone speech shows as low pitch variation.
Unstable vocal folds show as jitter and shimmer.
Incomplete closure of the folds shows as breathiness.
We measure that with HNR and CPPS.
Disturbed timing shows in the pause statistics.
Imprecise articulation shows in the spectral dynamics.
We wrote all seventy-four measurements ourselves.
Nothing was downloaded ready-made.
\arstart
\begin{Arabic}
لم نختر القياسات عشوائياً.
فكل قياس يعكس شيئاً مادياً في الحنجرة أو الفم.
الكلام الرتيب يظهر في قلة تغيّر النبرة.
وعدم استقرار الحبال الصوتية يظهر في الارتجاف والتذبذب.
وعدم اكتمال إغلاق الحبال يظهر في همس الصوت.
ونقيس ذلك بنسبة التوافقيات إلى الضجيج وبروز القمة الكبسترالية.
واضطراب الإيقاع يظهر في إحصاءات التوقّفات.
وضعف دقة مخارج الحروف يظهر في ديناميكيات الطيف.
وقد كتبنا القياسات الأربعة والسبعين كلها بأنفسنا.
لم نحمّل شيئاً جاهزاً.
\end{Arabic}
\clearpage

% ---------------- Slide 7 ----------------
\slidehead{Slide 7 --- We chose not to clean the sound}{target 0:55}
Here is one engineering decision, told honestly.
We did not remove noise from the recordings.
That sounds wrong. It is not.
Noise removal works by smoothing irregularities.
But irregularity is exactly our medical evidence.
Jitter and shimmer are irregularities.
If we clean the sound, a sick voice looks healthy.
So we protect the evidence and skip denoising.
\arstart
\begin{Arabic}
هذا قرار هندسي واحد، نرويه بصدق.
لم نُزل الضجيج من التسجيلات.
يبدو ذلك خطأ. لكنه ليس كذلك.
فإزالة الضجيج تعمل بتنعيم عدم الانتظام.
وعدم الانتظام هو بالضبط دليلنا الطبي.
فالارتجاف والتذبذب هما عدم انتظام.
لو نظّفنا الصوت، لبدا الصوت المريض سليماً.
لذلك نحمي الدليل ونتجنّب إزالة الضجيج.
\end{Arabic}
\clearpage

% ---------------- Slide 8 ----------------
\slidehead{Slide 8 --- A student with the exam questions}{target 1:25}
Now the heart of the project.
Imagine a student who gets the exam questions before the exam.
He scores one hundred percent.
Does that prove he understands? No.
The same danger exists here, in a hidden form.
Each person gave us two recordings.
Suppose one goes to training and one goes to testing.
The model can then simply recognise the person's voice.
It wins without detecting any disease.
This is called data leakage.
Our program checks for it in every split.
If it ever happens, the program stops itself.
\arstart
\begin{Arabic}
والآن قلب المشروع.
تخيّلوا طالباً حصل على أسئلة الامتحان قبل الامتحان.
سيحصل على مئة بالمئة.
هل يثبت ذلك أنه فاهم؟ لا.
والخطر نفسه موجود هنا، بشكل خفي.
كل شخص أعطانا تسجيلين.
افترضوا أن أحدهما ذهب للتدريب والآخر للاختبار.
عندها يستطيع النموذج أن يتعرّف على صوت الشخص فقط.
فيفوز من دون أن يكشف أي مرض.
وهذا اسمه تسريب البيانات.
وبرنامجنا يفحص ذلك في كل تقسيم.
وإذا حدث يوماً، يوقف البرنامج نفسه.
\end{Arabic}
\clearpage

% ---------------- Slide 9 ----------------
\slidehead{Slide 9 --- Validation design decides the number}{target 1:25}
Look at this number: \textcolor{honest}{0.909}.
It looks like an excellent result.
[Press for the next step.]
Now the truth.
This number comes from a deliberately wrong split.
Chunks of the same recording were on both sides.
The honest split gives 0.775.
Same data. Same model. Only the split changed.
The difference is memorisation, not detection.
This is why many papers report ninety-five percent or more.
Orange means inflated. Blue means honest.
Remember these two colours.
\arstart
\begin{Arabic}
انظروا إلى هذا الرقم: \textenglish{0.909}.
يبدو نتيجة ممتازة.
[اضغط للخطوة التالية.]
والآن الحقيقة.
هذا الرقم من تقسيم خاطئ متعمَّد.
كانت مقاطع التسجيل نفسه على الجانبين.
أما التقسيم الأمين فيعطي \textenglish{0.775}.
البيانات نفسها. والنموذج نفسه. تغيّر التقسيم فقط.
والفرق هو الحفظ، لا الكشف.
ولهذا تعلن أبحاث كثيرة خمسة وتسعين بالمئة أو أكثر.
البرتقالي يعني مضخَّماً. والأزرق يعني أميناً.
تذكّروا هذين اللونين.
\end{Arabic}
\clearpage

% ---------------- Slide 10 ----------------
\slidehead{Slide 10 --- The rules were fixed before the experiments}{target 1:00}
We wrote our selection rules before running any experiment.
Same splits for every variant.
A complex variant must win by more than two points.
Any result above ninety-five percent stops the study.
And all nineteen configurations are reported.
Now the punchline.
Our best score was 0.840, from a model ensemble.
We rejected it.
It won by only 0.018.
That is smaller than the random noise between runs.
Rejecting our best number is a strength, not a loss.
\arstart
\begin{Arabic}
كتبنا قواعد الاختيار قبل تشغيل أي تجربة.
التقسيمات نفسها لكل المتغيّرات.
والمتغيّر الأعقد يجب أن يتفوق بأكثر من نقطتين.
وأي نتيجة فوق خمسة وتسعين بالمئة توقف الدراسة.
وكل التكوينات التسعة عشر تُذكر في التقرير.
والآن بيت القصيد.
أفضل نتيجة كانت \textenglish{0.840}، من نموذج تجميعي.
ورفضناها.
لأنها تفوقت بمقدار \textenglish{0.018} فقط.
وهذا أصغر من الضجيج العشوائي بين التشغيلات.
ورفض أفضل رقم عندنا قوة، وليس خسارة.
\end{Arabic}
\clearpage

% ---------------- Slide 11 ----------------
\slidehead{Slide 11 --- Judged on people it has never heard}{target 1:05}
Here are the final numbers.
Balanced accuracy is \textcolor{honest}{0.822}, plus or minus 0.032.
Sensitivity is 0.771. Specificity is 0.873.
The area under the curve is 0.864.
Every number is measured on unseen people.
Why balanced accuracy?
A model that always says ``healthy'' gets 0.575 accuracy.
But it finds zero patients.
Balanced accuracy gives that model only 0.500.
So it cannot hide.
\arstart
\begin{Arabic}
هذه هي الأرقام النهائية.
الدقة المتوازنة \textenglish{0.822}، زائد أو ناقص \textenglish{0.032}.
والحساسية \textenglish{0.771}. والنوعية \textenglish{0.873}.
والمساحة تحت المنحنى \textenglish{0.864}.
وكل رقم مقيس على أشخاص لم يرهم النموذج.
لماذا الدقة المتوازنة؟
نموذج يقول دائماً «سليم» يحصل على دقة \textenglish{0.575}.
لكنه لا يجد أي مريض.
والدقة المتوازنة تعطيه \textenglish{0.500} فقط.
فلا يستطيع أن يختبئ.
\end{Arabic}
\clearpage

% ---------------- Slide 12 ----------------
\slidehead{Slide 12 --- It rediscovered a known clinical sign}{target 0:55}
We asked the model which measurement mattered most.
The answer: pitch range.
Pitch range is the acoustic form of monotone speech.
Doctors have described monotone speech for decades.
So a system with no medical knowledge rediscovered a clinical sign.
This is why we chose an explainable model.
A deep network could not answer this question.
\arstart
\begin{Arabic}
سألنا النموذج: أي قياس كان الأهم؟
الجواب: مدى النبرة.
ومدى النبرة هو الشكل الصوتي للكلام الرتيب.
والأطباء يصفون الكلام الرتيب منذ عقود.
أي أن نظاماً بلا معرفة طبية أعاد اكتشاف علامة سريرية.
ولهذا اخترنا نموذجاً قابلاً للتفسير.
فالشبكة العميقة لا تستطيع الإجابة عن هذا السؤال.
\end{Arabic}
\clearpage

% ---------------- Slide 13 ----------------
\slidehead{Slide 13 --- Italy: a frozen model, a trap, and a puzzle}{target 1:15}
We froze the model. No retraining. No adjustment.
Then we tested it on an Italian corpus. Sixty-one speakers.
First, inspection caught a trap.
Every high-quality file belonged to a patient.
The model could win by detecting the microphone.
So we limited all audio to sixteen kilohertz.
The result: area under the curve 0.701.
Balanced accuracy 0.629.
Performance dropped, but did not collapse.
Now a puzzle. Look at the young healthy speakers.
They scored higher than the real patients. Why?
[Pause. Let them think.]
Medically, this is impossible.
So the score follows the recording, not the disease.
That is the strongest argument for caution.
\arstart
\begin{Arabic}
جمّدنا النموذج. بلا إعادة تدريب. وبلا تعديل.
ثم اختبرناه على قاعدة إيطالية. واحد وستون متحدثاً.
أولاً، كشف الفحص فخّاً.
كل ملف عالي الجودة كان لمريض.
كان يمكن للنموذج أن يفوز بكشف الميكروفون.
لذلك حددنا كل الصوتيات عند ستة عشر كيلوهرتز.
والنتيجة: المساحة تحت المنحنى \textenglish{0.701}.
والدقة المتوازنة \textenglish{0.629}.
انخفض الأداء، لكنه لم ينهَر.
والآن لغز. انظروا إلى الشباب الأصحّاء.
حصلوا على درجات أعلى من المرضى الحقيقيين. لماذا؟
[توقّف. دعهم يفكرون.]
طبياً، هذا مستحيل.
إذن الدرجة تتبع التسجيل، لا المرض.
وهذه أقوى حجة للحذر.
\end{Arabic}
\clearpage

% ---------------- Slide 14 ----------------
\slidehead{Slide 14 --- A system that refuses to guess}{target 0:45}
This is the working prototype.
The disclaimer appears before any input.
Scores between 0.35 and 0.65 give ``inconclusive''.
Remember the Italian test.
Every elderly healthy speaker landed in that band.
One hundred percent of them.
The system refused to mislabel them.
Refusing to answer is a feature, not a failure.
\arstart
\begin{Arabic}
هذا هو النموذج الأولي وهو يعمل.
يظهر التنويه قبل أي إدخال.
والدرجات بين \textenglish{0.35} و\textenglish{0.65} تعطي «غير حاسم».
تذكّروا الاختبار الإيطالي.
كل متحدث سليم كبير في السن وقع في هذا النطاق.
مئة بالمئة منهم.
رفض النظام أن يصنّفهم تصنيفاً خاطئاً.
ورفض الإجابة ميزة، وليس فشلاً.
\end{Arabic}
\clearpage

% ---------------- Demo slide (optional) ----------------
\slidehead{Optional demo slide}{no time budget --- about 1:30 if used}
Use this slide only if time allows.
Say: I can show the system live.
Run the command line on a corpus recording.
The analysis takes about one minute.
Talk over the wait: the pipeline is extracting features now.
If anything misbehaves, stay calm.
Say: the screenshots on the previous slide show the same result.
Then continue to the closing slide.
\arstart
\begin{Arabic}
استخدم هذه الشريحة فقط إذا سمح الوقت.
قل: أستطيع عرض النظام مباشرة.
شغّل سطر الأوامر على تسجيل من قاعدة البيانات.
يستغرق التحليل نحو دقيقة واحدة.
تكلّم أثناء الانتظار: الأنبوب يستخرج الخصائص الآن.
وإذا تعطّل شيء، فابقَ هادئاً.
قل: لقطات الشاشة في الشريحة السابقة تعرض النتيجة نفسها.
ثم انتقل إلى شريحة الختام.
\end{Arabic}
\clearpage

% ---------------- Slide 15 ----------------
\slidehead{Slide 15 --- What this is, and is not}{target 0:35}
Let me close with what this system is not.
It does not diagnose. Many conditions change a voice.
It was trained on thirty-seven people, one language.
Now my single main message.
On small voice datasets, validation design moves the number
more than model engineering does.
An honest \textcolor{honest}{0.822} is worth more than an inflated one.
Thank you.
\arstart
\begin{Arabic}
أختم بما ليس عليه هذا النظام.
إنه لا يشخّص. فحالات كثيرة تغيّر الصوت.
وقد تدرّب على سبعة وثلاثين شخصاً وبلغة واحدة.
والآن رسالتي الأساسية الوحيدة.
في قواعد البيانات الصوتية الصغيرة، يحرّك تصميمُ التحقّق الرقمَ
أكثر مما تحرّكه هندسة النموذج.
ورقم أمين قدره \textenglish{0.822} أثمن من رقم مضخَّم.
شكراً لكم.
\end{Arabic}
\clearpage

% ================= BACKUP SLIDES (Q&A) =================
\slidehead{Backup 1 --- Why 82\% and not 95\%?}{Q\&A only}
Published high numbers usually split recordings, not people.
We measured that effect: 0.909 leaky against 0.775 honest.
Our 0.822 is measured on people the model never heard.
\arstart
\begin{Arabic}
الأرقام العالية المنشورة تُقسّم التسجيلات عادة، لا الأشخاص.
وقد قسنا هذا الأثر: \textenglish{0.909} بالتسريب مقابل \textenglish{0.775} بأمانة.
ورقمنا \textenglish{0.822} مقيس على أشخاص لم يسمعهم النموذج.
\end{Arabic}
\clearpage

\slidehead{Backup 2 --- Why not deep learning?}{Q\&A only}
Thirty-seven subjects is far too few for deep networks.
Our model can name the measurements it used.
Our small neural network collapsed in the confounder test.
\arstart
\begin{Arabic}
سبعة وثلاثون شخصاً عدد أقل بكثير مما تحتاجه الشبكات العميقة.
ونموذجنا يستطيع تسمية القياسات التي اعتمد عليها.
وشبكتنا العصبية الصغيرة انهارت في اختبار العامل المُربك.
\end{Arabic}
\clearpage

\slidehead{Backup 3 --- Why balanced accuracy?}{Q\&A only}
Our classes are not equal in size.
``Always healthy'' scores 0.575 accuracy and finds nobody.
Balanced accuracy exposes that model with 0.500.
\arstart
\begin{Arabic}
مجموعتانا غير متساويتين في الحجم.
ونموذج «الجميع أصحّاء» يحقق دقة \textenglish{0.575} ولا يجد أحداً.
والدقة المتوازنة تفضحه بإعطائه \textenglish{0.500}.
\end{Arabic}
\clearpage

\slidehead{Backup 4 --- Why reject the 0.840 ensemble?}{Q\&A only}
The margin of 0.02 was fixed before any experiment.
The ensemble won by 0.018, less than seed noise of 0.03.
Choosing the most complex of tied options fits the test, not the problem.
\arstart
\begin{Arabic}
هامش \textenglish{0.02} ثُبّت قبل أي تجربة.
والنموذج التجميعي تفوّق بمقدار \textenglish{0.018}، أقل من ضجيج البذور البالغ \textenglish{0.03}.
واختيار الأعقد من خيارات متعادلة يلائم الاختبار، لا المشكلة.
\end{Arabic}
\clearpage

\slidehead{Backup 5 --- Can it diagnose Parkinson's?}{Q\&A only}
No, and it is not designed to.
Colds, ageing, smoking, and fatigue also change a voice.
It reports similarity to research groups; clinicians make diagnoses.
\arstart
\begin{Arabic}
لا، وليس مصمّماً لذلك.
فالزكام والتقدّم في العمر والتدخين والإرهاق تغيّر الصوت أيضاً.
هو يُبلّغ عن التشابه مع مجموعات بحثية؛ والأطباء هم من يشخّصون.
\end{Arabic}
\clearpage

\slidehead{Backup 6 --- Is 37 subjects too few?}{Q\&A only}
Yes, and we say so openly in the report.
We compensated: grouped validation, three seeds, spread reported.
And we tested externally on sixty-one more speakers.
\arstart
\begin{Arabic}
نعم، ونقول ذلك صراحة في التقرير.
وقد عوّضنا: تحقّق مجمَّع، وثلاث بذور، وذكرنا التشتّت.
واختبرنا خارجياً على واحد وستين متحدثاً إضافياً.
\end{Arabic}
\clearpage

\slidehead{Backup 7 --- Is it detecting male vs female voices?}{Q\&A only}
Absolute pitch does encode the speaker's sex.
So we removed those four features and re-ran everything.
Our model held: 0.768 to 0.780. The neural network collapsed.
\arstart
\begin{Arabic}
النبرة المطلقة تعكس فعلاً جنس المتحدث.
لذلك حذفنا تلك الخصائص الأربع وأعدنا كل التقييم.
وصمد نموذجنا: من \textenglish{0.768} إلى \textenglish{0.780}. وانهارت الشبكة العصبية.
\end{Arabic}
\clearpage

\slidehead{Backup 8 --- The hardest technical problem?}{Q\&A only}
Finding the traps that would flatter our results.
In the Italian corpus, every 44.1 kHz file was a patient.
We caught it at inspection and band-limited everything to 16 kHz.
\arstart
\begin{Arabic}
أصعب شيء كان إيجاد الفخاخ التي كانت ستجمّل نتائجنا.
في القاعدة الإيطالية، كل ملف بمعدّل \textenglish{44.1} كيلوهرتز كان لمريض.
كشفنا ذلك عند الفحص وحددنا نطاق كل الملفات عند \textenglish{16} كيلوهرتز.
\end{Arabic}

\end{document}
